\documentclass[10pt,a4paper]{article}
\usepackage[margin=2.2cm]{geometry}
\usepackage{fontspec}
\setmonofont{DejaVu Sans Mono}[Scale=0.84]
\newfontfamily\unifont{DejaVu Serif}[Scale=0.9]
\usepackage{amsmath,amssymb,mathtools}
\setcounter{MaxMatrixCols}{20}
\usepackage{longtable,booktabs,array}
\usepackage{fancyvrb}
\usepackage[dvipsnames]{xcolor}
\usepackage[colorlinks=true,linkcolor=NavyBlue,urlcolor=NavyBlue]{hyperref}
\usepackage{microtype}
\usepackage{needspace}
\usepackage{tocloft}\setlength{\cftbeforesecskip}{0.45em}
\setlength{\parindent}{0pt}\setlength{\parskip}{0.55em}
\setlength{\emergencystretch}{8em}
\renewcommand{\arraystretch}{1.15}
\allowdisplaybreaks
\begin{document}
\title{The self-coupled dirac16complex in gravitational fields: provenance\\[0.4em]\large Sources, decisions, derivations and the verification record of the nine-channel self-coupling, the Fierz matrix, the covariant nonlinear field equation and the Kohn–Sham theory of dirac16complex in the three epochs — sections 15–16 of the notebook 7+1.nb}
\author{Repository: github.com/once-ere/7plus1 — prepared for Nash\\ Notebook: 7+1.nb (Wolfram Language 15.0.1)\\ Document generated by tools/make\_docs.py on 2026-10-27}
\date{}
\maketitle
\begin{abstract}
This record states what was asked, what was used, every decision taken with the equation behind it, how each result was derived and how it was checked: the 24 checks of sections 15–16 of 7+1.nb (24 True), the 29 boolean checks of the independent Python program tools/ref\_ks\_gravity.py (29 True), the comparison of every number printed by both programs, and the SHA-256 hashes of the files.
\end{abstract}
\tableofcontents
\Needspace{12\baselineskip}\section*{1. The task}\addcontentsline{toc}{section}{1. The task}
This stage of the project asked (summarized; the original wording is the author’s private working material):

\begin{enumerate}
\item a GENERALIZED self-coupling of the 16-component field dirac16complex, beyond the four contact couplings (scalar, pseudoscalar, vector, axial) of the earlier Kohn–Sham theory of the notebook;
\item the Kohn–Sham theory of dirac16complex in GRAVITATIONAL fields, with the canonical spin connection coupling gravity to the field, in the primordial epoch (near $x_0=0$), the modern epoch and the far future;
\item everything on the 7+1 spacetime of this project — one time coordinate $x_0$, seven space coordinates $x_1,\dots,x_7$ — with the ideas (not the 4+4 metric) of the repository github.com/once-ere/Dirac\_claude;
\item for every documentation file a companion provenance file and a symbol file, in Markdown and LaTeX with PDFs; no document refers the reader to another one; every equation used is written in the document that uses it;
\item and (27 October 2026) to make clear in all documentation, examples and provenance records that the xAct list $\{7,1,0\}$ declares seven space-like dimensions and one time-like dimension, while $\mathrm{diag}(+1,-1,\dots,-1)$ is simply the squares of the octonion basis (section 1.1).
\end{enumerate}
\Needspace{8\baselineskip}\subsection*{1.1 The signature: the author’s explanation of 27 October 2026}\addcontentsline{toc}{subsection}{1.1 The signature: the author’s explanation of 27 October 2026}
\textbf{Source.} On 27 October 2026 the author of the project explained that the xAct list $\{7,1,0\}$ declares seven space-like dimensions and one time-like dimension, while $\mathrm{diag}(+1,-1,\dots,-1)$ is simply the squares of the octonion basis, and asked that this be made clear in all documentation, examples and provenance records. \textbf{What it corrected.} The version of 7+1.nb pushed in the commits 52d8729 to dc2f178 of github.com/once-ere/7plus1 declared the xAct metric of section 12 with \texttt{DefMetric[\{1, 7, 0\}, …]}, reading the list as the signs of $\eta$; since then section 12 declares \texttt{DefMetric[\{7, 1, 0\}, g71[-a, -b], CD, …]}, exactly as the author’s notebook Einstein-Lovelock-7+1-Nash.nb does. \textbf{What was added.} Checks in 7+1.nb — section 4.7: $\eta$ is the squares of the octonion basis, $e_Ae_A=\eta_{AA}e_0$ and $\mathrm{Re}(e_Ae_B)=\eta_{AB}$; section 12.2: \texttt{SignatureOfMetric[g71]} $=\{7,1,0\}$ with \texttt{SignDetOfMetric[g71]} $=-1$, $\mathrm{Re}(e_Ae_B)=\eta_{AB}$, and $\mathrm{Re}(x\,x)=\eta_{AB}x^Ax^B$; section 17: the flat metric $\mathrm{diag}(+1,-1,-1,-1)$ of the world brane is the squares of the basis $e_0,e_1,e_2,e_3$ of the quaternion subalgebra — and the same statements in the Python programs \texttt{tools/ref71.py}, \texttt{tools/ref\_gravity.py} and \texttt{tools/ref\_sm.py}; the explanation itself in the text of 7+1.nb (sections 1.2, 1.4, 4.2, 4.7, 7, 11, 12.2, 13, 15.1, 17) and in every document, symbol file and provenance file of the repository. \textbf{What did not change.} No result. Compared output by output with the backup below, every output of 7+1.nb is the same, except the declared list itself, the new checks and the check counts (and, as in every run, the name of the private context and the measured times); in particular every component of every Lovelock tensor is unchanged. The components follow from the component matrix of the metric, and the sign of the determinant, \texttt{SignDetOfMetric}, is $-1$ for both lists. \textbf{Backup.} The evaluated notebook as it was before is kept in the workspace as \texttt{backups/7+1\_2026-10-27\_before\_signature.nb} (the folder backups/ is excluded from the repository by .gitignore), 1034401 bytes, SHA-256 \texttt{060942488caf4f45120ffd37be0123e96992a31cd7fdbde09ee4dc4ced4f5d3a}; \texttt{tools/verify\_7p1\_nb.py} compares every input, text and section cell of 7+1.nb with it and finds every cell present, each one unchanged or changed by this explanation, and it compares every output in the same way.

\textbf{The explanation, in full.} The 7+1 spacetime has one time-like dimension ($x_0$, the direction of the octonion unit $e_0=1$) and seven space-like dimensions ($x_1,x_2,x_3$ ordinary space, $x_4,\dots,x_7$ auxiliary space). Two different objects describe this, and neither may be replaced by the other. (i) The tensor package xAct, used in section 12 of 7+1.nb exactly as in the author’s notebook Einstein-Lovelock-7+1-Nash.nb, declares the metric with \texttt{DefMetric[\{7, 1, 0\}, g71[-a, -b], CD, …]}: the list $\{7,1,0\}$ declares seven space-like dimensions, one time-like dimension and no degenerate dimension — the numbers of entries $+1$ and $-1$ of $\eta_{71}=\mathrm{diag}(-1,+1,\dots,+1)$ (\texttt{η71}), the same spacetime written in the sign convention of 4+4.nb (space-like $=+1$, time-like $=-1$). (ii) The matrix $\eta=-\eta_{71}=\mathrm{diag}(+1,-1,\dots,-1)$ (\texttt{η17}) is simply the squares of the octonion basis: $e_0^2=+1$ for the unit $e_0=1$ (time) and $e_i^2=-1$ for the seven imaginary units $e_1,\dots,e_7$ (space); more generally it is the real part of the product of two basis units, and it gives the quadratic form of the octonions:
\begin{gather*}
\eta_{AB}=\mathrm{Re}(e_Ae_B),\qquad e_0^2=+1,\qquad e_i^2=-1\ \ (i=1,\dots,7),\\ \mathrm{Re}(x\,x)=x_0^2-x_1^2-\cdots-x_7^2=\eta_{AB}\,x^Ax^B\qquad (x=x^Ae_A\in\mathbb O).
\end{gather*}
$\eta$ gives the signs of the Clifford relation $\Gamma_A\Gamma_B+\Gamma_B\Gamma_A=2\,\eta_{AB}\,1_{16}$ of the real 16×16 Dirac matrices: time squares to $+1_{16}$ and each of the seven space directions to $-1_{16}$, exactly like the octonion units (in the notation $Cl(p,q)$, with $p$ generators squaring to $+1$ and $q$ to $-1$, the algebra $Cl(1,7)\cong\mathbb R(16)$). It also gives the signs of the metric components: $g_{00}>0$ along the time-like direction, negative squares along the space-like directions. Every component computed by 7+1.nb follows from the component matrix of the metric; the declared list fixes for xAct the sign of the determinant of the metric, \texttt{SignDetOfMetric[g71]} $=-1$, which is the same in both sign conventions, since $\det \eta=\det \eta_{71}=-1$. 7+1.nb checks \texttt{SignatureOfMetric[g71]} $=\{7,1,0\}$, \texttt{SignDetOfMetric[g71]} $=-1$, $\mathrm{Re}(e_Ae_B)=\eta_{AB}$ and $\mathrm{Re}(x\,x)=\eta_{AB}x^Ax^B$.

\Needspace{12\baselineskip}\section*{2. Inputs}\addcontentsline{toc}{section}{2. Inputs}
\Needspace{8\baselineskip}\subsection*{2.1 Objects of the earlier sections of 7+1.nb}\addcontentsline{toc}{subsection}{2.1 Objects of the earlier sections of 7+1.nb}
Used unchanged: the real Dirac matrices \texttt{Γ17[A]}, \texttt{Γ17up[A]}, the metric \texttt{η17}, the chirality \texttt{χ17} $=i\,\Gamma_0\Gamma_1\cdots\Gamma_7$, the chiral basis \texttt{EL}, \texttt{ER}, the Hamiltonian matrices \texttt{αD[j]}, \texttt{βD}, the curvature engine and the frame data of the cosmological metric (\texttt{fdCos}), the closed forms of the Lovelock tensors of the cosmological metric (\texttt{ECos}, with the rewriting \texttt{toH} in terms of $H_a,H_b,H_a',H_b'$), and the Hartree couplings and saturation state of the uniform matter of section 8.7 of the notebook:
\begin{align*}
\Gamma_0&=\begin{pmatrix}0&1_8\\1_8&0\end{pmatrix},\quad \Gamma_i=\begin{pmatrix}\tau^O_i&0\\0&-\tau^O_i\end{pmatrix},\quad \chi=i\,\Gamma_0\Gamma_1\cdots\Gamma_7, \\
E_L&=\tfrac1{\sqrt2}\begin{pmatrix}1_8\\-i1_8\end{pmatrix},\quad E_R=\tfrac1{\sqrt2}\begin{pmatrix}-i1_8\\1_8\end{pmatrix}, \\
G_S&=-100000,\quad G_V=10000\ \ (m=1), \\
k_{F,\rm sat}&=0.75663615070149289469\ldots\ \ \text{(where } E/A \text{ is minimal at } a=b=1).
\end{align*}
\Needspace{8\baselineskip}\subsection*{2.2 Software}\addcontentsline{toc}{subsection}{2.2 Software}
\begin{itemize}
\item Wolfram Language 15.0.1 (front end WolframNB.exe), with the standard package NumericalDifferentialEquationAnalysis (its function \texttt{GaussianQuadratureWeights}) for the Gauss–Legendre nodes.
\item Python 3.13.17 with numpy 2.5.3, scipy 1.18.1, sympy 1.14.0 (the independent program).
\item Textbook results, each derived and checked here: Fierz rearrangement (M. Fierz, Z. Phys. 104 (1937) 553), Hohenberg–Kohn and Kohn–Sham (Phys. Rev. 136 (1964) B864; 140 (1965) A1133), relativistic mean-field matter (J. D. Walecka, Ann. Phys. 83 (1974) 491).
\end{itemize}
\Needspace{12\baselineskip}\section*{3. Decisions}\addcontentsline{toc}{section}{3. Decisions}
{\footnotesize
\begin{longtable}{>{\raggedright\arraybackslash}p{\dimexpr 0.0550\linewidth-2\tabcolsep\relax}>{\raggedright\arraybackslash}p{\dimexpr 0.4750\linewidth-2\tabcolsep\relax}>{\raggedright\arraybackslash}p{\dimexpr 0.4700\linewidth-2\tabcolsep\relax}}
\toprule
\textbf{\#} & \textbf{decision} & \textbf{reason / equation} \\
\midrule\endfirsthead
\toprule
\textbf{\#} & \textbf{decision} & \textbf{reason / equation} \\
\midrule\endhead
\endfoot
\bottomrule\endlastfoot
1 & basis of the 256 matrices: $\Theta^K=i^{\sigma_r}\Gamma^{A_1}\cdots\Gamma^{A_r}$, $\sigma_r=\frac{r(r-1)}2\bmod2$ & $(\Gamma_0\Gamma^{A_1}\cdots\Gamma^{A_r})^\dagger=(-1)^{r(r-1)/2}\,\Gamma_0\Gamma^{A_1}\cdots\Gamma^{A_r}$ (reversing $r$ distinct anticommuting factors); the phase $i^{\sigma_r}$ makes $\Gamma_0\Theta^K$ Hermitian, so every bilinear is real (section 4) \\
2 & the self-coupling is organized by rank, one coupling $g_r$ per rank & each rank is a Lorentz-invariant class (the generators $\frac14[\Gamma^A,\Gamma^B]$ map products of rank $r$ to rank $r$) \\
3 & the four couplings of section 8 are identified as $g_0=g_S$, $g_1=g_V$, $g_7=-g_A$, $g_8=-g_P$ & $I_8=-(\bar\psi i\chi\psi)^2$ and $I_7=-(\bar\psi\Gamma^A\chi\psi)(\bar\psi\Gamma_A\chi\psi)$ \\
4 & exchange is folded into the direct terms by the Fierz matrix & $E_x=-\sum_{r,s}\frac{g_r}2F_{rs}D_s$, hence $G=(1-F^{\mathsf T})g$ \\
5 & uniform matter: spin-saturated, no Dirac sea, Hartree–Fock couplings $G_S=G_0$, $G_V=G_1$ & only $D_0=\rho_S^2$ and $D_1=\rho^2$ are non-zero for $R=(\rho_S1+\rho\Gamma_0)/16$ \\
6 & comoving Fermi ball: the comoving momenta are conserved by a homogeneous expansion & the modes $e^{i(\vec p\cdot\vec x+\vec q\cdot\vec y)}$ keep $\vec p,\vec q$; the physical momenta are $\vec p/a$, $\vec q/b$ \\
7 & the comoving Fermi momentum is the flat-space saturation value & the matter starts (at $a=b=1$) in its ground state at zero pressure \\
8 & adiabatic (instantaneous) Kohn–Sham state & valid when $H_a,H_b\ll$ orbital energies; the pair creation of the epochs measures the departures \\
9 & Gauss–Legendre quadrature with 96 × 64 nodes in $(k,\theta)$ & smooth integrands; the flat limit reproduces the closed forms of section 8.7 to $10^{-11}$ \\
10 & the 00 equation is imposed on the initial data only & it is a constraint without second derivatives; its preservation is a check of the whole system \\
11 & NDSolve with explicit Runge–Kutta of order 8 and InterpolationOrder → All & without it the interpolation between steps (order 3) showed an apparent constraint violation of $2\times10^{-6}$ \\*
\end{longtable}}
\Needspace{12\baselineskip}\section*{4. Derivations}\addcontentsline{toc}{section}{4. Derivations}
\textbf{Reality of the bilinears.} $\Gamma_0\Gamma^{A\dagger}\Gamma_0=\Gamma^A$, so $(\Gamma_0\Gamma^{A_1}\cdots\Gamma^{A_r})^\dagger=\Gamma^{A_r\dagger}\cdots\Gamma^{A_1\dagger}\Gamma_0=\Gamma_0\Gamma^{A_r}\cdots\Gamma^{A_1}=(-1)^{r(r-1)/2}\Gamma_0\Gamma^{A_1}\cdots\Gamma^{A_r}$ (distinct indices anticommute). Multiplying by $i^{\sigma_r}$ with $\sigma_r=\frac{r(r-1)}2\bmod2$ compensates the sign: $(i^{\sigma_r})^*(-1)^{r(r-1)/2}=i^{\sigma_r}$.

\textbf{Fierz matrix.} Completeness, $\sum_K(\Theta_K)_{cd}(\Theta^K)_{ab}=16\delta_{ad}\delta_{bc}$, gives $\sum_{K\in r}(a\Theta_Kb)(c\Theta^Kd)=\sum_JC_{rJ}(a\Theta_Jd)(c\Theta^Jb)$ with $C_{rJ}=\frac1{256}\sum_{K\in r}\mathrm{tr}(\Theta_K\Theta_J\Theta^K\Theta^J)$. Moving $\Gamma^{A}$ through a product of $s$ distinct $\Gamma$’s gives $(-1)^s$ if $A$ is not among them and $(-1)^{s-1}$ if it is; for $K$ of rank $r$ sharing $l$ indices with $J$ the sign is $(-1)^{rs-l}$, and there are $\binom sl\binom{8-s}{r-l}$ such $K$:
\begin{equation*}
F_{rs}=\frac{16}{256}\sum_l(-1)^{rs-l}\binom{s}{l}\binom{8-s}{r-l}=\frac1{16}\sum_l(-1)^{rs-l}\binom{s}{l}\binom{8-s}{r-l}.
\end{equation*}
\textbf{Densities of uniform matter.} A 16-component field has 8 positive-energy states per momentum, so $\rho=8\int\frac{d^7k}{(2\pi)^7}$ over the Fermi ball (comoving, divided by the comoving volume factor $a^3b^4$). With $k_a=k\cos\theta$, $k_b=k\sin\theta$:
\begin{gather*}
d^3k_a\,d^4k_b=(4\pi)(2\pi^2)\,k^6\cos^2\theta\sin^3\theta\,dk\,d\theta,\qquad \frac{8\cdot8\pi^3}{(2\pi)^7}=\frac1{2\pi^4},\\ \int_0^{\pi/2}\cos^2\theta\sin^3\theta\,d\theta=\frac2{15},\qquad \rho=\frac{k_F^7}{105\pi^4a^3b^4}.
\end{gather*}
The kinetic pressures are the momentum flux per direction: three ordinary directions share $k_a^2/a^2$, four auxiliary directions share $k_b^2/b^2$, which gives $p_a=\langle k^2\cos^2\theta/(3a^2\omega)\rangle+\Pi$ and $p_b=\langle k^2\sin^2\theta/(4b^2\omega)\rangle+\Pi$; the interaction pressure is $\Pi=\rho\,\partial\varepsilon_{\rm int}/\partial\rho-\varepsilon_{\rm int}$ of the mean-field energy $\varepsilon_{\rm int}=-\frac{(M-m)^2}{2G_S}+\frac{G_V}2\rho^2$ at fixed $M$.

\textbf{Conservation.} $\varepsilon(a,b,M)$ is stationary in $M$ (that is the gap equation), so $\varepsilon'=\partial_a\varepsilon\,a'+\partial_b\varepsilon\,b'$; differentiating the comoving integrals ($\rho\propto a^{-3}b^{-4}$, $\partial_a\omega=-k^2\cos^2\theta/(a^3\omega)$) gives
\begin{equation*}
\varepsilon'=-3H_a(\varepsilon+p_a)-4H_b(\varepsilon+p_b).
\end{equation*}
\textbf{Constraint preservation.} The Lovelock tensors are divergence-free for every metric, $\nabla_\mu E^{(k)\mu}{}_\nu=0$; for the cosmological metric the $\nu=0$ component reads $\frac{d}{dx_0}E^{(k)0}{}_0+3H_a(E^{(k)0}{}_0-E^{(k)1}{}_1)+4H_b(E^{(k)0}{}_0-E^{(k)4}{}_4)=0$. With the 11 and 44 equations imposed and the conservation law above, the 00 equation therefore holds at all times if it holds initially.

\Needspace{12\baselineskip}\section*{5. Verification}\addcontentsline{toc}{section}{5. Verification}
\textbf{Evaluation of 7+1.nb}: Wolfram 15.0.1 for Microsoft Windows (64-bit) (July 2, 2026); finished Tue 27 Oct 2026 19:31:35; message cells 0. The checks of sections 15–16:

{\footnotesize
\begin{longtable}{>{\raggedright\arraybackslash}p{\dimexpr 0.0550\linewidth-2\tabcolsep\relax}>{\raggedright\arraybackslash}p{\dimexpr 0.7450\linewidth-2\tabcolsep\relax}>{\raggedright\arraybackslash}p{\dimexpr 0.0800\linewidth-2\tabcolsep\relax}>{\raggedright\arraybackslash}p{\dimexpr 0.1200\linewidth-2\tabcolsep\relax}}
\toprule
\textbf{\#} & \textbf{check} & \textbf{result} & \textbf{section} \\
\midrule\endfirsthead
\toprule
\textbf{\#} & \textbf{check} & \textbf{result} & \textbf{section} \\
\midrule\endhead
\endfoot
\bottomrule\endlastfoot
1 & \texttt{256 multi-indices in 9 ranks with 1, 8, 28, 56, 70, 56, 28, 8, 1 members} & True & 15.1 \\
2 & \texttt{every bilinear psibar Theta\textasciicircum{}K psi is real: Gamma17[0].Theta\textasciicircum{}K Hermitian for all 256 K} & True & 15.1 \\
3 & \texttt{Theta\_K.Theta\textasciicircum{}K == ID16 and Tr[Theta\_K.Theta\textasciicircum{}J] == 16 delta\_KJ} & True & 15.1 \\
4 & \texttt{Hodge duality: Theta\textasciicircum{}K.omega17 == (phase) Theta\textasciicircum{}(complement of K), |phase| = 1} & True & 15.1 \\
5 & \texttt{the couplings of section 8 are four of the nine: I\_0 == (psibar psi)\textasciicircum{}2, I\_1 == (psibar Gamma\textasciicircum{}A psi)(psibar Gamma\_A psi), I\_8 == -(psibar i chi17 psi)\textasciicircum{}2, I\_7 == -(axial invariant)} & True & 15.1 \\
6 & \texttt{local Lorentz invariance: I\_r(S psi, psibar S\textasciicircum{}-1) == I\_r(psi, psibar) for r = 0..8 (random finite Lorentz transformation), and S\textasciicircum{}-1 == Gamma17[0].S\textasciicircum{}dagger.Gamma17[0]} & True & 15.1 \\
7 & \texttt{Fierz matrix: closed form == (1/256) Sum\_\{K in r\} Tr[Theta\_K.Theta\_J.Theta\textasciicircum{}K.Theta\textasciicircum{}J] for the first and last J of every rank s} & True & 15.2 \\
8 & \texttt{F.F == IdentityMatrix[9]; eigenvalues: five +1 and four -1} & True & 15.2 \\
9 & \texttt{Fierz identity for commuting spinors (random): Sum\_\{K in r\} (a Theta\_K b)(c Theta\textasciicircum{}K d) == Sum\_s F\_rs Sum\_\{J in s\} (a Theta\_J d)(c Theta\textasciicircum{}J b), r = 0..8} & True & 15.2 \\
10 & \texttt{Hartree-Fock energy E\_H + E\_x == Sum\_s (G\_s/2) D\_s with G = (1 - Transpose[F]) g (random density matrix and couplings)} & True & 15.2 \\
11 & \texttt{uniform density matrix (rhoS ID16 + rho Gamma17[0])/16: D\_0 == rhoS\textasciicircum{}2, D\_1 == rho\textasciicircum{}2, D\_2 .. D\_8 == 0} & True & 15.2 \\
12 & \texttt{section 8 recovered: with g = (gS, gV, 0, 0, 0, 0, 0, -gA, -gP) the exchange energy of uniform matter is that of section 8.7} & True & 15.2 \\
13 & \texttt{the Hartree-Fock couplings of section 8.7: gS = -100000, gV = 10000 give G\_0 = -98750, G\_1 = 20000} & True & 15.2 \\
14 & \texttt{nonlinear field equation: dL\_self/d psibar\_a == -Sum\_r g\_r Sum\_K (psibar Theta\_K psi) (Theta\textasciicircum{}K psi)\_a} & True & 15.3 \\
15 & \texttt{chi17 commutes with every generator (Gamma17up[A].Gamma17up[B] - Gamma17up[B].Gamma17up[A])/4, hence with every Omega\_mu} & True & 15.4 \\
16 & \texttt{curved chiral equations: E\_L\textasciicircum{}dag.Gamma17[0].(i Gamma\textasciicircum{}mu D\_mu - m) psi == (i d0 + i(3Ha/2 + 2Hb) + tauO.d) phi\_L - m phi\_R, and the same for E\_R with -tauO.d} & True & 15.4 \\
17 & \texttt{flat limit a = b = 1, kF = kFsat: m* == 0.3480973401... and E/A - m == -0.0283214279... (section 8.7), to 10\textasciicircum{}-11} & True & 16.1 \\
18 & \texttt{flat limit: pa == pb == rho eps\_F - eps (Gibbs-Duhem), which is 0 at the saturation density} & True & 16.1 \\
19 & \texttt{energy-momentum conservation of the Kohn-Sham matter: eps' + 3Ha(eps + pa) + 4Hb(eps + pb) == 0 (a = 1 + 0.3x + 0.05x\textasciicircum{}2, b = 1/(1 + 0.2x), x0 = 0.7; relative 10\textasciicircum{}-8)} & True & 16.1 \\
20 & \texttt{primordial: rho x0 == kFsat\textasciicircum{}7/(105 Pi\textasciicircum{}4) at every x0, and m* decreases monotonically to below 0.05 at x0 = 10\textasciicircum{}-4} & True & 16.2 \\
21 & \texttt{modern (b = 1): eps a\textasciicircum{}3 tends to a constant and pa/eps -> 0 (pressureless in ordinary space)} & True & 16.2 \\
22 & \texttt{far future (de Sitter, H = 0.1): rho == rho0 e\textasciicircum{}\{-3 H x0\} and 1 - m* decays like rho} & True & 16.2 \\
23 & \texttt{the 11 and 44 Einstein-Lovelock equations are linear in (Ha', Hb'); the 00 equation contains neither (constraint)} & True & 16.3 \\
24 & \texttt{self-consistent evolution: the constraint Sum alpha\_k E\textasciicircum{}(k)0\_0 - Lambda == kappa eps is preserved to 10\textasciicircum{}-7 (relative) for Einstein and for Einstein-Lovelock (alpha2 = 20, alpha3 = 500)} & True & 16.3 \\*
\end{longtable}}
\textbf{The independent Python program} tools/ref\_ks\_gravity.py — its boolean checks:

{\footnotesize
\begin{longtable}{>{\raggedright\arraybackslash}p{\dimexpr 0.0550\linewidth-2\tabcolsep\relax}>{\raggedright\arraybackslash}p{\dimexpr 0.8250\linewidth-2\tabcolsep\relax}>{\raggedright\arraybackslash}p{\dimexpr 0.1200\linewidth-2\tabcolsep\relax}}
\toprule
\textbf{\#} & \textbf{check} & \textbf{result} \\
\midrule\endfirsthead
\toprule
\textbf{\#} & \textbf{check} & \textbf{result} \\
\midrule\endhead
\endfoot
\bottomrule\endlastfoot
1 & \texttt{256 products in 9 ranks: 1, 8, 28, 56, 70, 56, 28, 8, 1} & True \\
2 & \texttt{every bilinear psibar Theta\textasciicircum{}K psi is real: Gamma17[0] Theta\textasciicircum{}K Hermitian (Theta\textasciicircum{}K = i\textasciicircum{}((r(r-1)/2) mod 2) Gamma\textasciicircum{}K)} & True \\
3 & \texttt{Theta\_K Theta\textasciicircum{}K = 1\_16 and tr(Theta\_K Theta\textasciicircum{}J) = 16 delta\_KJ} & True \\
4 & \texttt{Fierz completeness: sum\_K (Theta\_K)\_cd (Theta\textasciicircum{}K)\_ab = 16 delta\_ad delta\_bc} & True \\
5 & \texttt{Hodge duality: Theta\textasciicircum{}K omega17 = +- Theta\textasciicircum{}(complement of K) (up to a phase)} & True \\
6 & \texttt{stage-2 channels: I\_0 = (psibar psi)\textasciicircum{}2, I\_1 = (psibar Gamma\textasciicircum{}A psi)(psibar Gamma\_A psi), I\_8 = -(psibar i chi psi)\textasciicircum{}2, I\_7 = -(axial)} & True \\
7 & \texttt{local Lorentz invariance: I\_r(S psi, psibar S\textasciicircum{}-1) = I\_r(psi, psibar), r = 0..8 (random finite boost-rotation)} & True \\
8 & \texttt{Fierz matrix F\_rs = (1/256) sum\_\{K in r\} tr(Theta\_K Theta\_J Theta\textasciicircum{}K Theta\textasciicircum{}J), independent of the representative J of rank s} & True \\
9 & \texttt{Fierz matrix closed form F\_rs = (1/16) sum\_l (-1)\textasciicircum{}(rs-l) C(s,l) C(8-s,r-l)} & True \\
10 & \texttt{F\textasciicircum{}2 = 1 (the Fierz transformation is an involution)} & True \\
11 & \texttt{Fierz identity for commuting spinors: sum\_\{K in r\} (a Theta\_K b)(c Theta\textasciicircum{}K d) = sum\_s F\_rs sum\_\{J in s\} (a Theta\_J d)(c Theta\textasciicircum{}J b)} & True \\
12 & \texttt{Hartree-Fock energy E\_H + E\_x = sum\_s (G\_s/2) D\_s with G = (1 - F\textasciicircum{}T) g (random density matrix, random couplings)} & True \\
13 & \texttt{uniform matter R = (rho\_S + rho Gamma17[0])/16: only D\_0 = rho\_S\textasciicircum{}2, D\_1 = rho\textasciicircum{}2; with g\_0 = gS, g\_1 = gV, g\_7 = -gA, g\_8 = -gP the exchange energy is that of stage 2 (section 8.7)} & True \\
14 & \texttt{uniform density matrix: D\_0 = rho\_S\textasciicircum{}2, D\_1 = rho\textasciicircum{}2, D\_2..D\_8 = 0} & True \\
15 & \texttt{stage-2 couplings gS = -100000, gV = 10000: G\_0 = -98750, G\_1 = 20000 (the HF couplings of section 8.7)} & True \\
16 & \texttt{nonlinear field equation: dL\_int/d psibar\_a = -sum\_r g\_r sum\_K (psibar Theta\_K psi)(Theta\textasciicircum{}K psi)\_a} & True \\
17 & \texttt{the chirality chi17 commutes with every generator Sigma\textasciicircum{}AB, hence with every spin connection Omega\_mu} & True \\
18 & \texttt{curved chiral equations: E\_L\textasciicircum{}dag Gamma0 (i Gamma\textasciicircum{}mu D\_mu - m) psi = (i d0 + i(3Ha/2 + 2Hb) + tauO.d/(a,b)) phi\_L - m phi\_R} & True \\
19 & \texttt{curved chiral equations: E\_R\textasciicircum{}dag Gamma0 (i Gamma\textasciicircum{}mu D\_mu - m) psi = (i d0 + i(3Ha/2 + 2Hb) - tauO.d/(a,b)) phi\_R - m phi\_L} & True \\
20 & \texttt{flat limit a = b = 1, kF = kF0 (stage-2 saturation): m* = 0.3480973401..., E/A - m = -0.0283214279...} & True \\
21 & \texttt{flat limit: isotropic pressure pa = pb = rho eps\_F - eps (Gibbs-Duhem, eps\_F = V0 + sqrt(kF\textasciicircum{}2 + m*\textasciicircum{}2))} & True \\
22 & \texttt{energy-momentum conservation eps' + 3Ha(eps + pa) + 4Hb(eps + pb) = 0 for the KS matter (arbitrary a, b; relative residual)} & True \\
23 & \texttt{primordial: rho(x0) x0 = kF0\textasciicircum{}7/(105 pi\textasciicircum{}4) (density grows like 1/x0)} & True \\
24 & \texttt{primordial: m* decreases toward 0 as x0 -> 0 (blue-shifted ordinary momenta, attractive scalar coupling)} & True \\
25 & \texttt{modern: eps a\textasciicircum{}3 -> constant and pa/eps -> 0 (dirac16complex behaves as pressureless matter in ordinary space)} & True \\
26 & \texttt{far future: rho = rho0 e\textasciicircum{}\{-3 H x0\} and 1 - m* -> 0 like rho} & True \\
27 & \texttt{the 00 closed forms of E\textasciicircum{}(1), E\textasciicircum{}(2), E\textasciicircum{}(3) (pair-set algorithm, tools/lovelock71.py) agree with the closed forms of section 12.6 of 7+1.nb} & True \\
28 & \texttt{Einstein-Lovelock equations are linear in (Ha', Hb') (quasi-linear second-order system)} & True \\
29 & \texttt{the constraint Sum alpha\_k E\textasciicircum{}(k)0\_0 - Lambda = kappa eps is preserved by the evolution (relative violation < 1e-7) in both cases} & True \\*
\end{longtable}}
\textbf{Numbers printed by both programs} (relative differences):

{\small
\begin{longtable}{>{\raggedright\arraybackslash}p{\dimexpr 0.6000\linewidth-2\tabcolsep\relax}>{\raggedright\arraybackslash}p{\dimexpr 0.1500\linewidth-2\tabcolsep\relax}>{\raggedright\arraybackslash}p{\dimexpr 0.2500\linewidth-2\tabcolsep\relax}}
\toprule
\textbf{quantity} & \textbf{values} & \textbf{max. relative difference} \\
\midrule\endfirsthead
\toprule
\textbf{quantity} & \textbf{values} & \textbf{max. relative difference} \\
\midrule\endhead
\endfoot
\bottomrule\endlastfoot
flat limit: $m^*$, $\rho$, $\varepsilon$ & 3 & $7.5\times10^{-15}$ \\
primordial $m^*(x_0)$ & 6 & $5.3\times10^{-15}$ \\
primordial $p_a/p_b$ & 3 & $2.3\times10^{-15}$ \\
modern $m^*(a)$ & 5 & $4.1\times10^{-15}$ \\
modern $p_b/\varepsilon$ at $a=100$ & 1 & $3.1\times10^{-15}$ \\
far future $m^*(x_0)$ & 5 & $4.1\times10^{-15}$ \\
Einstein run: $H_b(0)$; $a,b,H_a,H_b$ at $x_0=20$ & 5 & $2.6\times10^{-12}$ \\
Einstein–Lovelock run: $H_b(0)$; $a,b,H_a,H_b$ at $x_0=20$ & 5 & $2.1\times10^{-12}$ \\*
\end{longtable}}
\textbf{Checks made while generating the documents:} 16 F printed by 7+1.nb == closed form: True.

\textbf{SHA-256 of the files} (at the time this document was generated):

\begin{Verbatim}[fontsize=\small,frame=single,framesep=2mm]
7+1.nb                             1314b1e09e6b6eefbd9ee5974e19d186a5e996907b2eb68ddb6f23aae85aa
    ...aa5
tools/make_7p1_nb.py               354a0210e60510499753c2cf91274850cedef502d6f705efc410c5a8bbae7
    ...b2d
tools/nb_part_gravity.py           c4caf2ad5c932e8f603341edade5f4b98fa784feb05906c051eab7c884cd6
    ...8f8
tools/nb_part_ks_gravity.py        a656c0ad040b4576b75b17e029545109e531c8e656ad99bfdb5cba946e967
    ...772
tools/nb_part_sm.py                43e92beed57453a4b57931de6d603ac6c218efba6e4005133b69b8054a975
    ...c5b
tools/nb_part_coupled.py           49c61aeb6affe14ebbcccb32d2069504b9fc035313ea2fabbd0a3392a9613
    ...3f5
tools/lovelock71.py                d93ecbd0f82adb972218d59dfe32999a4af944945ca9758607f3a633e1af9
    ...5a1
tools/ref_gravity.py               f07cf8adaff585af89c87d19532f1aaa7db606ad88ae90483b2d1f41341ed
    ...f36
tools/ref_ks_gravity.py            b6488ef24b1c16aacd8ff6cf6a5a1ba3f316e082449ad97480022d1a26383
    ...27f
tools/ref_sm.py                    f3f80806893f42bbf166e1f1de896ffd611dbfe847685da8e3ee367e9cde6
    ...aba
tools/ref_coupled.py               695dbd9375753653472d9c37502f0067c4a64fe01a64ba8c05e8b2c5da85a
    ...a58
tools/ref71.py                     1a7a06eb0add3f459611776dc9c474a5bb959d947a9db949db4612b1ab943
    ...4e0
tools/ref44.py                     ab100c3ebcc055fbe59c132a02f80ecb17630e934e99d59f88e919ec8064f
    ...af5
\end{Verbatim}
\textbf{Output of tools/ref\_ks\_gravity.py (independent checks of sections 15-16; exit code 0)} (complete):

\begin{Verbatim}[fontsize=\small,frame=single,framesep=2mm]
256 products in 9 ranks: 1, 8, 28, 56, 70, 56, 28, 8, 1 : True
every bilinear psibar Theta^K psi is real: Gamma17[0] Theta^K Hermitian (Theta^K =
    ...i^((r(r-1)/2) mod 2) Gamma^K) : True
Theta_K Theta^K = 1_16 and tr(Theta_K Theta^J) = 16 delta_KJ : True
Fierz completeness: sum_K (Theta_K)_cd (Theta^K)_ab = 16 delta_ad delta_bc : True
Hodge duality: Theta^K omega17 = +- Theta^(complement of K) (up to a phase) : True
stage-2 channels: I_0 = (psibar psi)^2, I_1 = (psibar Gamma^A psi)(psibar Gamma_A psi), I_8 =
    ...-(psibar i chi psi)^2, I_7 = -(axial) : True
local Lorentz invariance: I_r(S psi, psibar S^-1) = I_r(psi, psibar), r = 0..8 (random finite
    ...boost-rotation) : True
Fierz matrix F_rs = (1/256) sum_{K in r} tr(Theta_K Theta_J Theta^K Theta^J), independent of
    ...the representative J of rank s : True
Fierz matrix closed form F_rs = (1/16) sum_l (-1)^(rs-l) C(s,l) C(8-s,r-l) : True
F^2 = 1 (the Fierz transformation is an involution) : True
eigenvalues of F: number of +1, number of -1 : (5, 4)
Fierz matrix (rows r, columns s) times 16 : [[1, 1, 1, 1, 1, 1, 1, 1, 1], [8, -6, 4, -2, 0, 2,
    ...-4, 6, -8], [28, 14, 4, -2, -4, -2, 4, 14, 28], [56, -14, -4, 6, 0, -6, 4, 14, -56],
    ...[70, 0, -10, 0, 6, 0, -10, 0, 70], [56, 14, -4, -6, 0, 6, 4, -14, -56], [28, -14, 4, 2,
    ...-4, 2, 4, -14, 28], [8, 6, 4, 2, 0, -2, -4, -6, -8], [1, -1, 1, -1, 1, -1, 1, -1, 1]]
Fierz identity for commuting spinors: sum_{K in r} (a Theta_K b)(c Theta^K d) = sum_s F_rs
    ...sum_{J in s} (a Theta_J d)(c Theta^J b) : True
Hartree-Fock energy E_H + E_x = sum_s (G_s/2) D_s with G = (1 - F^T) g (random density matrix,
    ...random couplings) : True
uniform matter R = (rho_S + rho Gamma17[0])/16: only D_0 = rho_S^2, D_1 = rho^2; with g_0 = gS,
    ...g_1 = gV, g_7 = -gA, g_8 = -gP the exchange energy is that of stage 2 (section 8.7) :
    ...True
uniform density matrix: D_0 = rho_S^2, D_1 = rho^2, D_2..D_8 = 0 : True
stage-2 couplings gS = -100000, gV = 10000: G_0 = -98750, G_1 = 20000 (the HF couplings of
    ...section 8.7) : True
effective couplings G_0, G_1 in terms of all nine g_r : (15*g0/16 - g1/2 - 7*g2/4 - 7*g3/2 -
    ...35*g4/8 - 7*g5/2 - 7*g6/4 - g7/2 - g8/16, -g0/16 + 11*g1/8 - 7*g2/8 + 7*g3/8 - 7*g5/8 +
    ...7*g6/8 - 3*g7/8 + g8/16)
example g = (-100000, 10000, 3000, 0, ..., 0): G_0, G_1 : (-104000, 17375)
nonlinear field equation: dL_int/d psibar_a = -sum_r g_r sum_K (psibar Theta_K psi)(Theta^K
    ...psi)_a : True
the chirality chi17 commutes with every generator Sigma^AB, hence with every spin connection
    ...Omega_mu : True
curved chiral equations: E_L^dag Gamma0 (i Gamma^mu D_mu - m) psi = (i d0 + i(3Ha/2 + 2Hb) +
    ...tauO.d/(a,b)) phi_L - m phi_R : True
curved chiral equations: E_R^dag Gamma0 (i Gamma^mu D_mu - m) psi = (i d0 + i(3Ha/2 + 2Hb) -
    ...tauO.d/(a,b)) phi_R - m phi_L : True
flat limit a = b = 1, kF = kF0 (stage-2 saturation): m* = 0.3480973401..., E/A - m =
    ...-0.0283214279... : True
flat limit: isotropic pressure pa = pb = rho eps_F - eps (Gibbs-Duhem, eps_F = V0 + sqrt(kF^2 +
    ...m*^2)) : True
NUM flat: m*, rho, eps, p : (0.3480973401462447, np.float64(1.3881002132316779e-05),
    ...np.float64(1.3487872330347983e-05), np.float64(-1.588186776101813e-20))
energy-momentum conservation eps' + 3Ha(eps + pa) + 4Hb(eps + pb) = 0 for the KS matter
    ...(arbitrary a, b; relative residual) : True
NUM conservation residual / eps' : -5.166108748755716e-12
primordial (Einstein-Kasner background, comoving kF = kF0 at x0 = 1): m*(x0) for x0 = 1e-4 ...
    ...10 : (0.003824179379430487, 0.010783989711922532, 0.030238519612042748,
    ...0.08844464116180238, 0.3480973401462447, 0.8917637493589158)
primordial: rho(x0) x0 = kF0^7/(105 pi^4) (density grows like 1/x0) : True
primordial: m* decreases toward 0 as x0 -> 0 (blue-shifted ordinary momenta, attractive scalar
    ...coupling) : True
primordial: anisotropic pressure pa/pb at x0 = 1e-3, 1e-1, 10 :
    ...(np.float64(1.0868869169150372), np.float64(1.6461737747098435),
    ...np.float64(-0.4036333858521348))
modern (static auxiliary space b = 1): m*(a) for a = 1, 2, 4, 10, 100 : (0.3480973401462447,
    ...0.853540809048541, 0.9807410717586909, 0.9987576007278518, 0.9999987563024887)
modern: eps a^3 -> constant and pa/eps -> 0 (dirac16complex behaves as pressureless matter in
    ...ordinary space) : True
modern: pb/eps at a = 100 (stress along the auxiliary space, balanced by the stabilizing
    ...stress) : 0.049773512087215956
far future (de Sitter a = e^{0.1 x0}, b = 1): m*(x0) for x0 = 0, 10, 20, 40, 80 -> m = 1 :
    ...(0.3480973401462447, 0.9395796910293632, 0.996923581335845, 0.9999923586141617,
    ...0.9999999999530483)
far future: rho = rho0 e^{-3 H x0} and 1 - m* -> 0 like rho : True
the 00 closed forms of E^(1), E^(2), E^(3) (pair-set algorithm, tools/lovelock71.py) agree with
    ...the closed forms of section 12.6 of 7+1.nb : True
Einstein-Lovelock equations are linear in (Ha', Hb') (quasi-linear second-order system) : True
Einstein + KS matter (kappa = 1000, Lambda = 0, Ha(0) = 0.05): Hb(0) from the constraint, (a,
    ...b, Ha, Hb) at x0 = 20, max relative constraint violation : (0.009143712444558836,
    ...(np.float64(1.5817386798469366), np.float64(1.1975585618423046),
    ...np.float64(0.011877999367582402), np.float64(0.007419107766685833)),
    ...np.float64(8.415965253866209e-12))
Einstein-Lovelock (alpha2 = 20, alpha3 = 500) + KS matter: Hb(0), (a, b, Ha, Hb) at x0 = 20,
    ...max relative constraint violation : (0.014980703942304356,
    ...(np.float64(1.4773514780943342), np.float64(1.276084053398673),
    ...np.float64(0.010175914492223711), np.float64(0.008488764693542736)),
    ...np.float64(4.387597609521105e-12))
the constraint Sum alpha_k E^(k)0_0 - Lambda = kappa eps is preserved by the evolution
    ...(relative violation < 1e-7) in both cases : True
boolean checks: 29   failed: 0
VERDICT: ALL CHECKS PASSED
\end{Verbatim}
\end{document}
